%%======================================================================
%%  MT-2: Literature Extraction & Convention Fixing
%%        for Modified Cosmology, H_0 and S_8 Tensions
%%
%%  This document collects formulas, data, and references needed to
%%  assess whether SCT's modified Friedmann equations can address the
%%  Hubble tension (H_0) and the growth-of-structure tension (S_8).
%%
%%  RESEARCH DOCUMENT: formulas extracted from cited papers with
%%  independent cross-checks.  No original derivations---all new results
%%  are in the downstream derivation document.
%%
%%  CENTRAL FINDING: SCT cannot resolve either tension through its
%%  FLRW-level nonlocal corrections, which are O(H^2/Lambda^2) << 1.
%%  This is a NEGATIVE result, honestly reported per research standard SR3.
%%======================================================================
\documentclass[11pt,a4paper]{article}

%% ---- packages -------------------------------------------------------
\usepackage[utf8]{inputenc}
\usepackage[T1]{fontenc}
\usepackage{lmodern}
\usepackage{amsmath,amssymb,amsthm,mathtools}
\usepackage{physics}
\usepackage[margin=2.5cm]{geometry}
\usepackage{booktabs}
\usepackage{enumitem}
\usepackage{hyperref}
\usepackage{cleveref}
\usepackage{xcolor}
\usepackage{fancyhdr}
\usepackage{tcolorbox}
\usepackage{longtable}
\usepackage{cancel}

%% ---- theorem-like environments --------------------------------------
\theoremstyle{plain}
\newtheorem{theorem}{Theorem}[section]
\newtheorem{proposition}[theorem]{Proposition}
\newtheorem{lemma}[theorem]{Lemma}
\newtheorem{corollary}[theorem]{Corollary}
\theoremstyle{definition}
\newtheorem{definition}[theorem]{Definition}
\newtheorem{convention}[theorem]{Convention}
\theoremstyle{remark}
\newtheorem{remark}[theorem]{Remark}

%% ---- custom commands (matching NT-4c conventions) --------------------
\providecommand{\Tr}{\operatorname{Tr}}
\providecommand{\tr}{\operatorname{tr}}
\newcommand{\Id}{\mathbb{1}}
\newcommand{\Ricsq}{R_{\mu\nu}R^{\mu\nu}}
\newcommand{\Rsq}{R^2}
\newcommand{\Weylsq}{C_{\mu\nu\rho\sigma}C^{\mu\nu\rho\sigma}}
\renewcommand{\dd}{\mathrm{d}}
\newcommand{\ie}{\textit{i.e.}}
\newcommand{\eg}{\textit{e.g.}}
\newcommand{\erfi}{\operatorname{erfi}}
\DeclareMathOperator{\sgn}{sgn}
\newcommand{\Lam}{\Lambda}

%% ---- annotation commands (from NT-1) --------------------------------
\newcommand{\fromlit}[1]{\textcolor{blue!70!black}{\textsf{[#1]}}}
\newcommand{\oursign}{\textcolor{teal!80!black}{\textsf{[our convention]}}}
\newcommand{\gap}{\textcolor{red!80!black}{\textsf{[GAP]}}}
\newcommand{\needscheck}{\textcolor{orange!80!black}{\textsf{[needs check]}}}

%% ---- verification annotations (from NT-4a) -------------------------
\newcommand{\verified}[1]{%
  \par\smallskip\noindent
  \colorbox{green!10}{\parbox{\dimexpr\linewidth-2\fboxsep}{%
    \textcolor{green!50!black}{\textbf{[Verified]}} #1}}%
  \par\smallskip}

\newcommand{\signcorrection}[1]{%
  \par\smallskip\noindent
  \colorbox{red!10}{\parbox{\dimexpr\linewidth-2\fboxsep}{%
    \textcolor{red!60!black}{\textbf{[Important note]}} #1}}%
  \par\smallskip}

\newcommand{\crosscheck}[1]{%
  \par\smallskip\noindent
  \colorbox{blue!10}{\parbox{\dimexpr\linewidth-2\fboxsep}{%
    \textcolor{blue!50!black}{\textbf{[Cross-check]}} #1}}%
  \par\smallskip}

\newcommand{\negativeresult}[1]{%
  \par\smallskip\noindent
  \colorbox{orange!10}{\parbox{\dimexpr\linewidth-2\fboxsep}{%
    \textcolor{orange!60!black}{\textbf{[Negative result]}} #1}}%
  \par\smallskip}

%% ---- header ---------------------------------------------------------
\pagestyle{fancy}
\fancyhf{}
\fancyhead[L]{\small SCT Theory --- Literature Review}
\fancyhead[R]{\small MT-2: $H_0$ and $S_8$ Tensions}
\fancyfoot[C]{\thepage}

\hypersetup{
  colorlinks=true,
  linkcolor=blue!60!black,
  citecolor=green!50!black,
  urlcolor=blue!50!black,
  pdfauthor={David Alfyorov},
  pdftitle={MT-2: Literature Review --- H0 and S8 Tensions in SCT},
  pdfsubject={Spectral Causal Theory --- Cosmological Tensions}
}

\title{MT-2: Literature Extraction\\[4pt]
The $H_0$ and $S_8$ Tensions in the Context of\\
Spectral Causal Theory\\[6pt]
\large Conventions, Observational Data, Comparison Models,\\
and an Honest Assessment of SCT's Reach}
\author{David Alfyorov}
\date{March 2026}

\begin{document}
\maketitle

\begin{center}
\small\textit{Credits: Aliaksandr Samatyia contributed research-assistance
and workflow support.}
\end{center}

\begin{abstract}
This document collects the observational data, theoretical formulas,
and comparison models needed to assess whether Spectral Causal Theory
(SCT) can address the Hubble constant tension ($H_0$) and the
growth-of-structure tension ($S_8$). Drawing on the modified Friedmann
equations derived in NT-4c and the full nonlinear field equations from
NT-4b, we identify the fundamental obstacle: the SCT nonlocal
corrections to the Friedmann equations scale as
$\mathcal{O}(\beta_R \cdot H^2/\Lambda^2)$, which for any
astrophysically viable $\Lambda$ yields corrections of order
$10^{-61}$ to $10^{-120}$---completely negligible at cosmological
scales. We compare this with the Maggiore group's ``RR model'' and the
Deser--Woodard model, which achieve cosmologically relevant
corrections through \emph{infrared} nonlocality ($\Box^{-1}$,
$\Box^{-2}$) with a mass scale $m \sim H_0$, fundamentally different
from SCT's \emph{ultraviolet} nonlocality ($F(\Box/\Lambda^2)$ with
$\Lambda \gg H_0$). We catalog three possible indirect channels
through which SCT might affect cosmological observables---modified
initial conditions, running of $G_N$, and the modified Poisson
equation---but note that none has been demonstrated to produce
observable effects. This is an honest negative result.
\end{abstract}

\tableofcontents
\newpage

%======================================================================
\section{Introduction: The Two Cosmological Tensions}
\label{sec:intro}
%======================================================================

\subsection{The $H_0$ Tension}

The Hubble constant $H_0$ parametrizes the present-day expansion rate
of the universe via $v = H_0 d$ for nearby galaxies. Two classes of
measurements yield discrepant values:

\begin{center}
\begin{tabular}{llll}
\toprule
\textbf{Measurement} & $H_0$ (km/s/Mpc) & \textbf{Method} & \textbf{Reference} \\
\midrule
Planck 2018 (CMB) & $67.36 \pm 0.54$ & CMB + $\Lambda$CDM & \cite{Planck18VI} \\
SH0ES 2022 & $73.04 \pm 1.04$ & Cepheids + SNe Ia & \cite{Riess22} \\
TRGB (Freedman+21) & $69.8 \pm 1.7$ & TRGB + SNe Ia & \cite{Freedman21} \\
H0LiCOW & $73.3^{+1.7}_{-1.8}$ & Strong lensing time delays & \cite{Wong20} \\
CCHP (Freedman+24) & $69.96 \pm 1.05$ & JWST TRGB + JAGB + Cepheids & \cite{Freedman24} \\
DESI DR1 + Planck & $67.97 \pm 0.38$ & BAO + CMB & \cite{DESI24} \\
\bottomrule
\end{tabular}
\end{center}

The tension between Planck and SH0ES exceeds $5\sigma$
\fromlit{Di~Valentino+ 2021}. The TRGB measurements yield an
intermediate value that is compatible with both at $\lesssim 2\sigma$,
though JWST observations have not resolved the discrepancy
\cite{Freedman24}.

\signcorrection{The tension significance depends on which local
measurement is used. Planck vs.\ SH0ES: $>5\sigma$. Planck vs.\ TRGB:
$\sim 1.5\sigma$. The ``$H_0$ problem'' is most acute for the
Cepheid-based distance ladder.}

\subsection{The $S_8$ Tension}

The parameter $S_8 \equiv \sigma_8\sqrt{\Omega_m/0.3}$ characterizes the
amplitude of matter clustering. The CMB predicts a higher value than
direct low-redshift measurements:

\begin{center}
\begin{tabular}{lll}
\toprule
\textbf{Measurement} & $S_8$ & \textbf{Reference} \\
\midrule
Planck 2018 (CMB) & $0.832 \pm 0.013$ & \cite{Planck18VI} \\
KiDS-1000 & $0.759^{+0.024}_{-0.021}$ & \cite{KiDS1000} \\
DES Y3 & $0.776 \pm 0.017$ & \cite{DESY3} \\
HSC Y3 & $0.766^{+0.033}_{-0.031}$ & \cite{HSC23} \\
ACT DR6 lensing & $0.818 \pm 0.022$ & \cite{ACT23} \\
\bottomrule
\end{tabular}
\end{center}

The tension is $\sim 2\text{--}3\sigma$ for individual experiments,
though it has slightly weakened with the most recent measurements
\fromlit{DES Y3, ACT DR6}. Modified gravity models that weaken
gravitational clustering at low redshifts can in principle reduce
$S_8$.

\subsection{Why These Tensions Matter for SCT}

The roadmap task MT-2 asks: can the modified Friedmann equations from
the spectral action address either tension? The answer depends on the
magnitude of the nonlocal corrections relative to the GR baseline.

%======================================================================
\section{SCT Modified Friedmann Equations (from NT-4c)}
\label{sec:sct-friedmann}
%======================================================================

\subsection{The Equations}

From NT-4c (verified, 117/117 tests), the modified first Friedmann
equation on a spatially flat FLRW background is \fromlit{NT4c,
Eq.~(4.1)}:
\begin{equation}\label{eq:mod-friedmann}
  H^2 = \frac{\kappa^2}{3}\,\rho
  - \frac{\kappa^2}{3}\,\beta_R\,
  \bigl[H_{00} + \Theta^{(R)}_{00}\bigr],
\end{equation}
where:
\begin{itemize}[nosep]
  \item $\kappa^2 = 8\pi G$,
  \item $\beta_R = \alpha_R(\xi)/(16\pi^2) = 2(\xi - 1/6)^2/(16\pi^2)$
    is the reduced $R^2$ coupling,
  \item $H_{00}$ is the 00-component of the $R^2$-sector tensor,
  \item $\Theta^{(R)}_{00}$ is the Barvinsky--Vilkovisky alpha-insertion.
\end{itemize}

\begin{convention}[Key SCT coefficients]\label{conv:coeffs}
\begin{align}
  \alpha_C &= \frac{13}{120} \quad (\text{Weyl}^2,\; \xi\text{-independent}), \\
  \alpha_R(\xi) &= 2\left(\xi - \frac{1}{6}\right)^2, \\
  \beta_R &= \frac{\alpha_R(\xi)}{16\pi^2}.
\end{align}
\end{convention}

\subsection{Critical Property: On FLRW, Only the $R^2$ Sector Survives}

Since $C_{\mu\nu\rho\sigma} = 0$ on any FLRW background (conformal
flatness), the \emph{entire} Weyl sector---$F_1$, the Bach tensor
$B_{\mu\nu}$, and $\Theta^{(C)}_{\mu\nu}$---drops out of the
background equations. Only the $R^2$ sector, proportional to
$\alpha_R(\xi)$, contributes corrections.

\subsection{Conformal Coupling Kills All Corrections}

At conformal coupling $\xi = 1/6$:
\begin{equation}
  \alpha_R(1/6) = 0 \;\;\Longrightarrow\;\; \beta_R = 0
  \;\;\Longrightarrow\;\;
  \text{Standard GR on FLRW, identically.}
\end{equation}

\subsection{Alpha-Insertion Equation of State}

The Barvinsky--Vilkovisky alpha-insertion $\Theta^{(R)}_{\mu\nu}$ has
$w_\Theta = +1$ (stiff matter) \fromlit{NT4c, Cor.~4.3}. This is a
kinetic contribution from time derivatives of curvature, not a vacuum
energy term.

\subsection{De Sitter Exactness}

De~Sitter ($H = H_0 = \text{const}$) is an exact solution: both
$H_{00}$ and $\Theta^{(R)}_{00}$ vanish identically when $\dot{R} = 0$.

\verified{All properties verified in NT-4c: 117/117 tests PASS, dual
derivation (Methods A and B1), 100-digit precision. EOS $w = +1$ at 5
random FLRW backgrounds. De~Sitter vanishing at 5 values of $H_0$.}

%======================================================================
\section{The Suppression Argument: Why SCT Cannot Help}
\label{sec:suppression}
%======================================================================

This section contains the central negative result of this literature
review.

\subsection{Order-of-Magnitude Estimate}

The correction to $H^2$ from the spectral action is proportional to
\begin{equation}\label{eq:suppression}
  \frac{\delta H^2}{H^2} \sim \beta_R \cdot \frac{H^2}{\Lambda^2}
  = \frac{2(\xi - 1/6)^2}{16\pi^2} \cdot \frac{H^2}{\Lambda^2}.
\end{equation}

For the present epoch, $H_0 \approx 2.2 \times 10^{-18}$~Hz
$= 1.44 \times 10^{-33}$~eV:

\begin{center}
\begin{tabular}{lccc}
\toprule
$\Lambda$ & $(H_0/\Lambda)^2$ & $\beta_R(\xi=0)$ & $\delta H^2/H^2$ \\
\midrule
$M_{\mathrm{Pl}} \approx 1.22 \times 10^{19}$~GeV
  & $\sim 10^{-122}$ & $\sim 3.5 \times 10^{-4}$ & $\sim 10^{-125}$ \\
$1$~GeV & $\sim 10^{-84}$ & $\sim 3.5 \times 10^{-4}$ & $\sim 10^{-87}$ \\
$1$~eV & $\sim 10^{-66}$ & $\sim 3.5 \times 10^{-4}$ & $\sim 10^{-69}$ \\
PPN-1 lower bound: $2.38 \times 10^{-3}$~eV
  & $\sim 3.6 \times 10^{-61}$ & $\sim 3.5 \times 10^{-4}$ & $\sim 1.3 \times 10^{-64}$ \\
\bottomrule
\end{tabular}
\end{center}

\negativeresult{Even at the PPN-1 lower bound $\Lambda \geq 2.38 \times
10^{-3}$~eV (the smallest $\Lambda$ compatible with solar system tests),
the correction to $H^2$ is of order $10^{-64}$. The $H_0$ tension
requires a correction of order $\delta H/H \sim (73 - 67)/67 \approx 9\%$,
\ie, $\delta H^2/H^2 \sim 18\%$. The SCT correction falls short by
\textbf{64 orders of magnitude} at a minimum.}

\subsection{Could a Very Low $\Lambda$ Help?}

To get $\delta H^2/H^2 \sim 0.1$, one would need
\begin{equation}
  \frac{H_0^2}{\Lambda^2} \sim \frac{0.1}{\beta_R} \sim 70
  \quad\Longrightarrow\quad
  \Lambda \sim \frac{H_0}{\sqrt{70}} \sim 10^{-34}\;\text{eV}.
\end{equation}

This is catastrophically below the PPN-1 lower bound
$\Lambda \geq 2.38 \times 10^{-3}$~eV by 31 orders of magnitude.
A cutoff $\Lambda \sim H_0$ would destroy all short-distance
predictions of the theory (the modified Newtonian potential, ghost
masses, and all UV structure).

\negativeresult{There is no consistent choice of $\Lambda$ that
simultaneously satisfies PPN-1 bounds and produces cosmologically
relevant corrections to $H(z)$.}

\subsection{Formal Reason: UV vs.\ IR Nonlocality}

The fundamental issue is that SCT generates \emph{ultraviolet}
nonlocality through the form factors $F_i(\Box/\Lambda^2)$, where
$\Lambda$ is the spectral cutoff (a UV scale). The corrections become
relevant when the curvature scale approaches $\Lambda$, \ie, in the
early universe near reheating or in strong-gravity regions.

Models that successfully address the $H_0$ tension (see
\cref{sec:comparison}) use \emph{infrared} nonlocality through
operators like $\Box^{-1}$ or $\Box^{-2}$, with a mass parameter
$m \sim H_0$ that is dynamically generated in the far infrared. These
are fundamentally different types of nonlocality:

\begin{center}
\begin{tabular}{lcc}
\toprule
\textbf{Property} & \textbf{SCT} & \textbf{Maggiore RR model} \\
\midrule
Nonlocal operator & $F(\Box/\Lambda^2)$ (entire function) & $\Box^{-2}$ (inverse) \\
Mass scale & $\Lambda \gg H_0$ (UV) & $m \sim H_0$ (IR) \\
Cosmological corrections & $\mathcal{O}(H^2/\Lambda^2) \ll 1$ & $\mathcal{O}(1)$ \\
Origin & One-loop effective action & Dynamical mass generation \\
$H_0$ tension reduction & None & $5\sigma \to 2\sigma$ \\
$w_{\mathrm{DE}}$ & Not applicable & $\approx -1.14$ (phantom) \\
\bottomrule
\end{tabular}
\end{center}

\crosscheck{Maggiore~\cite{Maggiore16} explicitly showed that the
perturbative loop corrections (which produce form factors of the SCT
type) cannot generate cosmologically relevant nonlocality. The
required mass scale is in the non-perturbative infrared regime. This
confirms that the SCT one-loop spectral action, by construction,
produces UV-type form factors that are irrelevant at cosmological
energies.}

%======================================================================
\section{Observational Data Catalog}
\label{sec:data}
%======================================================================

For completeness and for use in downstream tasks (LT-3c), we catalog
the relevant datasets.

\subsection{$H_0$ Measurements}

\begin{longtable}{llll}
\toprule
\textbf{Probe} & $H_0$ (km/s/Mpc) & \textbf{Type} & \textbf{Reference} \\
\midrule
\endfirsthead
\toprule
\textbf{Probe} & $H_0$ (km/s/Mpc) & \textbf{Type} & \textbf{Reference} \\
\midrule
\endhead
Planck 2018 TT,TE,EE+lowE+lensing & $67.36 \pm 0.54$ & Early (CMB) & \cite{Planck18VI} \\
SH0ES (Riess+ 2022) & $73.04 \pm 1.04$ & Late (Cepheids) & \cite{Riess22} \\
TRGB (Freedman+ 2021) & $69.8 \pm 1.7$ & Late (TRGB) & \cite{Freedman21} \\
CCHP (Freedman+ 2024) & $69.96 \pm 1.05$ & Late (JWST) & \cite{Freedman24} \\
H0LiCOW (Wong+ 2020) & $73.3^{+1.7}_{-1.8}$ & Late (lensing) & \cite{Wong20} \\
Megamaser (Pesce+ 2020) & $73.9 \pm 3.0$ & Late (direct) & \cite{Pesce20} \\
DESI DR1 + Planck & $67.97 \pm 0.38$ & Early+BAO & \cite{DESI24} \\
ACT DR4 + WMAP & $67.6 \pm 1.1$ & Early (CMB) & \cite{ACT21} \\
SPT-3G 2018 & $68.3 \pm 1.5$ & Early (CMB) & \cite{SPT23} \\
\bottomrule
\end{longtable}

\subsection{$S_8$ Measurements}

\begin{longtable}{llll}
\toprule
\textbf{Survey} & $S_8$ & \textbf{Method} & \textbf{Reference} \\
\midrule
\endfirsthead
\toprule
\textbf{Survey} & $S_8$ & \textbf{Method} & \textbf{Reference} \\
\midrule
\endhead
Planck 2018 & $0.832 \pm 0.013$ & CMB primary & \cite{Planck18VI} \\
KiDS-1000 & $0.759^{+0.024}_{-0.021}$ & Cosmic shear & \cite{KiDS1000} \\
DES Y3 & $0.776 \pm 0.017$ & $3\times 2$pt & \cite{DESY3} \\
HSC Y3 & $0.766^{+0.033}_{-0.031}$ & Cosmic shear & \cite{HSC23} \\
ACT DR6 lensing & $0.818 \pm 0.022$ & CMB lensing & \cite{ACT23} \\
Planck 2018 lensing & $0.811 \pm 0.019$ & CMB lensing & \cite{Planck18VIII} \\
\bottomrule
\end{longtable}

\subsection{BAO Data (DESI DR1)}

\begin{center}
\begin{tabular}{lccc}
\toprule
\textbf{Tracer} & $z_{\mathrm{eff}}$ & $D_V/r_d$ & \textbf{Reference} \\
\midrule
BGS & 0.295 & $7.93 \pm 0.15$ & \cite{DESI24} \\
LRG1 & 0.510 & $13.62 \pm 0.25$ & \cite{DESI24} \\
LRG2 & 0.706 & $16.85 \pm 0.32$ & \cite{DESI24} \\
LRG3+ELG1 & 0.930 & $21.71 \pm 0.28$ & \cite{DESI24} \\
ELG2 & 1.317 & $27.79 \pm 0.69$ & \cite{DESI24} \\
Ly$\alpha$ QSO & 2.330 & $39.71 \pm 0.94$ & \cite{DESI24} \\
\bottomrule
\end{tabular}
\end{center}

\subsection{Type Ia Supernovae}

\begin{itemize}[nosep]
  \item \textbf{Pantheon+} (Brout+ 2022): 1701 SNe Ia, $0.001 < z < 2.26$,
    with SH0ES Cepheid calibration \cite{Brout22}.
  \item \textbf{Union3} (Rubin+ 2023): updated compilation, consistent
    with Pantheon+ \cite{Rubin23}.
  \item \textbf{DES-SN5YR} (DES 2024): 1635 photometrically classified
    SNe Ia, $0.1 < z < 1.3$ \cite{DES5YR}.
\end{itemize}

%======================================================================
\section{Comparison Models: Nonlocal Gravity That Can Help}
\label{sec:comparison}
%======================================================================

\subsection{Maggiore's RR Model}

The ``RR model'' \cite{MaggioreMancarella14, Belgacem18} modifies the
gravitational effective action by adding an infrared nonlocal term:
\begin{equation}\label{eq:RR-model}
  S_{\mathrm{RR}} = \frac{M_{\mathrm{Pl}}^2}{2}
  \int \dd^4 x\,\sqrt{-g}\,
  \left[R - \frac{m^2}{6}\,R\,\Box^{-2}\,R\right],
\end{equation}
where $m \sim H_0$ is a \emph{dynamically generated} mass scale. Key
properties:

\begin{enumerate}[nosep]
  \item \textbf{Dark energy without $\Lambda$:} The nonlocal term
    generates an effective dark energy with phantom equation of state
    $w_{\mathrm{DE}}(z=0) \approx -1.14$ \cite{Belgacem18}.
  \item \textbf{$H_0$ tension reduction:} Bayesian parameter estimation
    yields $H_0 \approx 69.5$~km/s/Mpc, reducing the tension with
    SH0ES to $\sim 2\sigma$ \cite{Belgacem18}.
  \item \textbf{Structure formation:} Modified growth rate
    $\gamma \approx 0.53$ vs.\ $\gamma \approx 0.55$ in $\Lambda$CDM
    \cite{Dirian14}. Effective Newton's constant modified by
    $\sim 9\%$ at $z = 0$ \cite{Dirian14}.
  \item \textbf{GW speed:} $c_T = c$ (passes GW170817) \cite{Belgacem18}.
  \item \textbf{Lunar Laser Ranging:} The RR model is ruled out by
    $|\dot{G}|/G$ constraints; the RT variant survives
    \cite{Belgacem19LLR}.
\end{enumerate}

\signcorrection{The RR model was later disfavored by Lunar Laser
Ranging constraints on $|\dot{G}|/G$ \cite{Belgacem19LLR}. The
surviving variant (``RT model'') modifies the conformal mode differently
and predicts $w_{\mathrm{DE}} \approx -1.04$, closer to $\Lambda$CDM
but still with distinctive GW luminosity distance predictions
\cite{Belgacem19GW}.}

\subsection{Deser--Woodard Model}

The Deser--Woodard model \cite{DeserWoodard07} introduces
$f(\Box^{-1} R)$ corrections:
\begin{equation}\label{eq:DW-model}
  S_{\mathrm{DW}} = \frac{M_{\mathrm{Pl}}^2}{2}
  \int \dd^4 x\,\sqrt{-g}\,
  R\bigl[1 + f(\Box^{-1} R)\bigr].
\end{equation}
This model also uses infrared nonlocality ($\Box^{-1}$) and can
generate accelerating cosmologies. However, it was ruled out by the
same $|\dot{G}|/G$ constraint that disfavors the RR model
\cite{Belgacem19LLR}.

\subsection{Why These Models Differ Fundamentally from SCT}

The models in \cref{eq:RR-model,eq:DW-model} use the \emph{inverse}
d'Alembertian ($\Box^{-1}$, $\Box^{-2}$), which is an infrared
operator: it amplifies long-wavelength modes. SCT uses the spectral
action form factors $F_i(\Box/\Lambda^2)$, which are \emph{entire
functions} of $\Box/\Lambda^2$ (verified in NT-2). These suppress
short-wavelength modes and become trivial at long wavelengths:
\begin{equation}
  F_i(z) \;\xrightarrow{z \to 0}\; F_i(0) + F_i'(0)\,z + \cdots
  \quad\text{(Taylor expansion, $z = -k^2/\Lambda^2$).}
\end{equation}
For $k \sim H_0 \ll \Lambda$, the argument $z$ is negligible and
$F_i \approx F_i(0) = \text{const}$, recovering the local limit.

Maggiore~\cite{Maggiore16} explicitly demonstrated that perturbative
loop corrections (which produce entire-function form factors of the
type used in SCT) cannot generate the required infrared nonlocality.
The mechanism must be non-perturbative, involving dynamical mass
generation for the conformal mode---a qualitatively different physics
from the spectral action.

%======================================================================
\section{Modified Gravity Approaches to the $S_8$ Tension}
\label{sec:S8-approaches}
%======================================================================

\subsection{General Mechanism}

Modified gravity can reduce $S_8$ by weakening the effective
gravitational coupling for matter perturbations at low redshifts. In
GR, the growth of linear perturbations $\delta = \delta\rho/\rho$
obeys:
\begin{equation}\label{eq:growth-GR}
  \ddot{\delta} + 2H\dot{\delta} - 4\pi G_{\mathrm{eff}}\,\rho\,\delta = 0.
\end{equation}
If $G_{\mathrm{eff}} < G_N$ at low $z$, the growth is suppressed and
$S_8$ decreases.

\subsection{Nonlocal Gravity and Growth Suppression}

In the Maggiore RR model, the effective gravitational constant for
perturbations is modified \cite{Dirian14}:
\begin{equation}
  \Psi(a;k) = \bigl[1 + \mu_s\,a^s\bigr]\,\Psi_{\mathrm{GR}}(a;k),
  \qquad \mu_s \approx 0.09,\; s \approx 2,
\end{equation}
leading to a $\sim 3\text{--}4\%$ deviation in the matter power spectrum
at $z = 0$. This is a scale-independent modification (anisotropic stress
is negligible) \cite{Dirian14}.

\subsection{What SCT Predicts for Growth}

From the NT-4a linearized field equations, the modified Poisson equation
in SCT is:
\begin{equation}\label{eq:poisson-sct}
  k^2\Phi = -4\pi G \cdot \frac{1}{\Pi_{\mathrm{s}}(k^2/\Lambda^2,\xi)}\,
  a^2\rho\,\delta,
\end{equation}
where $\Pi_{\mathrm{s}}(z,\xi) = 1 + 6(\xi - 1/6)^2\,z\,\hat{F}_2(z,\xi)$
\fromlit{NT-4a, Eq.~(5.12)}.

For sub-horizon modes relevant to structure formation ($k \sim 0.01
\text{--}1\;\text{Mpc}^{-1}$), the argument is:
\begin{equation}
  z = \frac{k^2}{\Lambda^2} \sim
  \frac{(0.1\;\text{Mpc}^{-1})^2}{(2.38 \times 10^{-3}\;\text{eV})^2}
  \sim 10^{-56},
\end{equation}
using the PPN-1 lower bound. For $\Lambda \sim M_{\mathrm{Pl}}$:
$z \sim 10^{-119}$.

Therefore $\Pi_{\mathrm{s}} \approx 1 + \mathcal{O}(10^{-56})$, and the
Poisson equation is indistinguishable from GR.

\negativeresult{The SCT modification to the Poisson equation is of
order $10^{-56}$ to $10^{-119}$ at the scales relevant for structure
formation. This cannot affect $S_8$ at any observable level.}

%======================================================================
\section{Possible Indirect Channels}
\label{sec:indirect}
%======================================================================

While the direct FLRW corrections and the growth-equation modifications
are negligible, there are three indirect channels through which SCT
might conceivably affect cosmological observables. None has been
demonstrated to work.

\subsection{Channel 1: Modified Initial Conditions from Early Universe}

Near the end of inflation or reheating ($T \sim 10^{13}$~GeV), the
Hubble parameter $H \sim 10^{13}$~GeV may approach $\Lambda$ for
sub-Planckian $\Lambda$ (a scenario considered in INF-1). In this
regime, the spectral corrections are $\mathcal{O}(1)$ and could modify
the primordial power spectrum.

However:
\begin{itemize}[nosep]
  \item INF-1 showed that the scalaron mass is $M = \sqrt{24\pi^2}\,M_{\mathrm{Pl}}
    \approx 15.4\,M_{\mathrm{Pl}}$ for SM-only content at $\xi = 0$---far too heavy
    for successful inflation \cite{INF1}.
  \item The conditional predictions (\ie, if $M$ is externally fixed) yield
    $n_s = 0.965$, $r = 3.5 \times 10^{-3}$, which are compatible with
    $\Lambda$CDM-based Planck analysis. There is no mechanism to shift $H_0$.
\end{itemize}

\textbf{Status:} \gap{} Requires a viable inflationary model within SCT (pending
resolution of the scalaron mass problem from INF-1) to assess whether
modified initial conditions could affect late-time parameters.

\subsection{Channel 2: Running of $G_N$}

The SCT propagator modification implies a scale-dependent effective
Newton's constant:
\begin{equation}\label{eq:G-running}
  G_{\mathrm{eff}}(k) = \frac{G_N}{\Pi_{\mathrm{TT}}(k^2/\Lambda^2)},
  \qquad
  \Pi_{\mathrm{TT}}(z) = 1 + \frac{13}{60}\,z\,\hat{F}_1(z).
\end{equation}

For cosmological momentum scales ($k \sim H_0$):
$z \sim 10^{-120}$ and $G_{\mathrm{eff}} = G_N(1 + \mathcal{O}(10^{-120}))$.

The running becomes significant only for $k \sim \Lambda$, which is
far above any cosmological scale.

\textbf{Status:} No effect at cosmological scales.

\subsection{Channel 3: Cosmological Constant from the Spectral Action}

The spectral action at zeroth order in curvature produces a cosmological
constant term:
\begin{equation}\label{eq:Lambda-CC}
  \Lambda_{\mathrm{CC}} \sim \frac{f_0\,\Lambda^4}{16\pi^2},
\end{equation}
where $f_0 = \int_0^\infty u\,f(u)\,\dd u$ is the zeroth moment of the
cutoff function. For $\Lambda \sim M_{\mathrm{Pl}}$, this gives
$\Lambda_{\mathrm{CC}} \sim M_{\mathrm{Pl}}^4$---the standard
cosmological constant problem.

The roadmap (MT-2 outcomes) asks whether nonlocal form-factor dressing
can suppress this bare $\Lambda^4$ contribution. The answer appears to
be no: the form factors $F_i(z)$ modify the curvature-squared sector,
not the zeroth-order cosmological constant. The $f_0$ term is a
\emph{tree-level} contribution that is not dressed by one-loop form
factors.

\textbf{Status:} \gap{} The cosmological constant problem in SCT
remains open, as in virtually all quantum gravity approaches.

%======================================================================
\section{Comparison Table: SCT vs.\ Models That Address Tensions}
\label{sec:comparison-table}
%======================================================================

\begin{longtable}{p{3.5cm}ccc}
\toprule
\textbf{Feature} & \textbf{SCT} & \textbf{Maggiore RT} & \textbf{$\Lambda$CDM} \\
\midrule
\endfirsthead
\toprule
\textbf{Feature} & \textbf{SCT} & \textbf{Maggiore RT} & \textbf{$\Lambda$CDM} \\
\midrule
\endhead
Free parameters (gravity) & 1 ($\Lambda$) + 1 ($\xi$) & 1 ($m$) & 0 \\
$w_{\mathrm{DE}}$ & Not applicable & $\approx -1.04$ & $-1$ (exact) \\
$H_0$ prediction & $= \Lambda$CDM & $\sim 69.5$ & $\sim 67.4$ \\
$H_0$ tension & $>5\sigma$ (unchanged) & $\sim 2\sigma$ & $>5\sigma$ \\
$S_8$ modification & None ($10^{-56}$) & $\sim 3\%$ & None \\
$c_T = c$? & Yes (structural) & Yes & Yes \\
$|\dot{G}|/G$ & Passes (SCT) & Passes (RT) & Passes \\
UV behavior & Entire functions & IR only & N/A \\
Cosmological const. & Open & Not addressed & Fine-tuned \\
\bottomrule
\end{longtable}

%======================================================================
\section{What SCT Can Do: Positive Cosmological Results}
\label{sec:positive}
%======================================================================

While SCT cannot address the $H_0$ or $S_8$ tensions, the FLRW
analysis yields several positive results that are reported honestly:

\begin{enumerate}
  \item \textbf{GW speed: $c_T = c$} (NT-4c, Theorem~5.1). This is a
    non-trivial survival test: GW170817 ruled out many modified gravity
    theories (Horndeski subclasses, massive gravity variants). SCT
    passes this test as a structural property of the action, not through
    fine-tuning.

  \item \textbf{De Sitter stability.} The current accelerating
    expansion (approximately de~Sitter) is a stable attractor of the
    modified equations. Perturbations decay exponentially.

  \item \textbf{No new cosmological degrees of freedom.} The theory does
    not introduce new light fields that could conflict with CMB
    observations (unlike some scalar-tensor theories).

  \item \textbf{Exact GR recovery at conformal coupling.} At $\xi = 1/6$,
    the theory is exactly GR on FLRW, providing a clean limit.

  \item \textbf{Consistency with all late-time observations.} Because
    the corrections are negligible, SCT is trivially consistent with
    all precision cosmological data---it predicts exactly $\Lambda$CDM
    at late times.
\end{enumerate}

%======================================================================
\section{Gaps Identified for Downstream Derivation}
\label{sec:gaps}
%======================================================================

\begin{center}
\begin{tabular}{clll}
\toprule
\textbf{\#} & \textbf{Gap} & \textbf{Impact} & \textbf{Resolution} \\
\midrule
G1 & Early-universe corrections & Unknown & Requires INF-1 completion \\
G2 & Cosmological constant & Open problem & Beyond one-loop \\
G3 & Non-perturbative IR mass & Not in SCT & Fundamental architecture \\
G4 & Causal set $\Lambda$ prediction & Not connected & MT-3 (Postulate 5) \\
G5 & $w_{\mathrm{eff}}(z)$ computation & Trivially $= -1$ & Not needed \\
G6 & Structure growth function & $= \Lambda$CDM & Not needed \\
G7 & DESI DR1 fit & $= \Lambda$CDM & Not needed \\
\bottomrule
\end{tabular}
\end{center}

Gaps G5--G7 are marked ``not needed'' because the direct computation
would yield $\Lambda$CDM predictions to machine precision, providing
no new information.

%======================================================================
\section{Conventions Summary}
\label{sec:conventions}
%======================================================================

\begin{convention}[Cosmological conventions]\label{conv:cosmo}
Throughout this analysis:
\begin{itemize}[nosep]
  \item Metric signature: $(-,+,+,+)$, FLRW: $\dd s^2 = -\dd t^2 + a(t)^2\,\dd\vec{x}^{\,2}$.
  \item $H \equiv \dot{a}/a$, $R = 6(\dot{H} + 2H^2)$.
  \item $\kappa^2 = 8\pi G$, $M_{\mathrm{Pl}} = 1/\sqrt{G} = 1.22 \times 10^{19}$~GeV.
  \item $S_8 \equiv \sigma_8\sqrt{\Omega_m/0.3}$.
  \item Critical density: $\rho_{\mathrm{cr}} = 3H_0^2/(8\pi G)$.
  \item $\Lambda$ is the spectral cutoff, \emph{not} the cosmological constant.
    The cosmological constant is written $\Lambda_{\mathrm{CC}}$ when needed.
\end{itemize}
\end{convention}

%======================================================================
\section{Summary and Verdict}
\label{sec:summary}
%======================================================================

\begin{tcolorbox}[colback=orange!5!white, colframe=orange!50!black,
  title={\textbf{MT-2 Literature Review: Central Finding}}]
\textbf{SCT cannot resolve the $H_0$ or $S_8$ tensions.}

The nonlocal corrections from the spectral action to the Friedmann
equations are of order
$\mathcal{O}(\beta_R \cdot H^2/\Lambda^2) \sim 10^{-64}$ to
$10^{-125}$, which is $64$--$125$ orders of magnitude below the
$\sim 10\%$ modification needed to address the $H_0$ tension. The
modification to the growth equation is similarly negligible.

This is because SCT generates \emph{ultraviolet} nonlocality (form
factors $F_i(\Box/\Lambda^2)$ with $\Lambda \gg H_0$), while
addressing cosmological tensions requires \emph{infrared} nonlocality
(operators like $\Box^{-1}$ with $m \sim H_0$).

At the same time, SCT passes all precision cosmological tests
($c_T = c$, de~Sitter stability, no new light fields) and is trivially
consistent with late-time data by reducing to $\Lambda$CDM.

This is an honest negative result (SR3), similar to INF-1's finding
that the scalaron mass is too heavy for inflation.
\end{tcolorbox}

\begin{tcolorbox}[colback=blue!5!white, colframe=blue!50!black,
  title={\textbf{Recommendation for Downstream Tasks}}]
\begin{enumerate}[nosep]
  \item \textbf{Derivation:} Confirm the suppression numerically at
    representative redshifts. Compute $\delta H^2/H^2$ vs.\ $z$ and
    $\Lambda$ to produce a definitive figure.
  \item \textbf{LT-3c:} Report that SCT predicts $\Lambda$CDM
    expansion history and growth function to machine precision.
  \item \textbf{COMP-1:} Contrast SCT's UV nonlocality with IR
    nonlocal models in the comparison table.
  \item \textbf{Do not attempt} a CLASS/cobaya fit: it would
    reproduce $\Lambda$CDM identically.
\end{enumerate}
\end{tcolorbox}

%======================================================================
\begin{thebibliography}{99}
%======================================================================

\bibitem{Planck18VI}
Planck Collaboration (N.~Aghanim \textit{et~al.}),
``Planck 2018 results. VI. Cosmological parameters,''
\emph{Astron.\ Astrophys.} \textbf{641} (2020) A6,
\href{https://arxiv.org/abs/1807.06209}{arXiv:1807.06209}.

\bibitem{Riess22}
A.~G.~Riess \textit{et~al.},
``A comprehensive measurement of the local value of the Hubble constant
with 1~km/s/Mpc uncertainty from the Hubble Space Telescope and the
SH0ES team,''
\emph{Astrophys.\ J.\ Lett.} \textbf{934} (2022) L7,
\href{https://arxiv.org/abs/2112.04510}{arXiv:2112.04510}.

\bibitem{Freedman21}
W.~L.~Freedman,
``Measurements of the Hubble constant: tensions in perspective,''
\emph{Astrophys.\ J.} \textbf{919} (2021) 16,
\href{https://arxiv.org/abs/2106.15656}{arXiv:2106.15656}.

\bibitem{Freedman24}
W.~L.~Freedman \textit{et~al.},
``Status report on the Chicago--Carnegie Hubble Program (CCHP):
three independent astrophysical determinations of the Hubble constant
using the James Webb Space Telescope,''
\href{https://arxiv.org/abs/2408.06153}{arXiv:2408.06153}.

\bibitem{Wong20}
K.~C.~Wong \textit{et~al.} (H0LiCOW),
``H0LiCOW---XIII. A 2.4\% measurement of $H_0$ from lensed quasars:
$5.3\sigma$ tension between early- and late-Universe probes,''
\emph{Mon.\ Not.\ Roy.\ Astron.\ Soc.} \textbf{498} (2020) 1420,
\href{https://arxiv.org/abs/1907.04869}{arXiv:1907.04869}.

\bibitem{Pesce20}
D.~W.~Pesce \textit{et~al.},
``The Megamaser Cosmology Project. XIII.\ Combined Hubble constant
constraints,''
\emph{Astrophys.\ J.\ Lett.} \textbf{891} (2020) L1,
\href{https://arxiv.org/abs/2001.09213}{arXiv:2001.09213}.

\bibitem{DESI24}
DESI Collaboration (A.~G.~Adame \textit{et~al.}),
``DESI 2024 VI: Cosmological constraints from the measurements of
baryon acoustic oscillations,''
\href{https://arxiv.org/abs/2404.03002}{arXiv:2404.03002}.

\bibitem{ACT21}
S.~Aiola \textit{et~al.} (ACT),
``The Atacama Cosmology Telescope: DR4 maps and cosmological parameters,''
\emph{JCAP} \textbf{12} (2020) 047,
\href{https://arxiv.org/abs/2007.07288}{arXiv:2007.07288}.

\bibitem{SPT23}
L.~Balkenhol \textit{et~al.} (SPT-3G),
``Measurement of the CMB temperature power spectrum and constraints on
cosmology from the SPT-3G 2018 TT, TE, and EE dataset,''
\emph{Phys.\ Rev.\ D} \textbf{108} (2023) 023510,
\href{https://arxiv.org/abs/2212.05642}{arXiv:2212.05642}.

\bibitem{KiDS1000}
M.~Asgari \textit{et~al.} (KiDS),
``KiDS-1000 cosmology: cosmic shear constraints on the amplitude of the
power spectrum,''
\emph{Astron.\ Astrophys.} \textbf{645} (2021) A104,
\href{https://arxiv.org/abs/2007.15633}{arXiv:2007.15633}.

\bibitem{DESY3}
DES Collaboration (T.~M.~C.~Abbott \textit{et~al.}),
``Dark Energy Survey Year 3 results: cosmological constraints from
galaxy clustering and weak lensing,''
\emph{Phys.\ Rev.\ D} \textbf{105} (2022) 023520,
\href{https://arxiv.org/abs/2105.13549}{arXiv:2105.13549}.

\bibitem{HSC23}
R.~Dalal \textit{et~al.} (HSC),
``Hyper Suprime-Cam Year 3 results: cosmology from cosmic shear
power spectra,''
\emph{Phys.\ Rev.\ D} \textbf{108} (2023) 123519,
\href{https://arxiv.org/abs/2304.00701}{arXiv:2304.00701}.

\bibitem{ACT23}
M.~S.~Madhavacheril \textit{et~al.} (ACT),
``The Atacama Cosmology Telescope: DR6 gravitational lensing map
and cosmological parameters,''
\emph{Astrophys.\ J.} \textbf{966} (2024) 138,
\href{https://arxiv.org/abs/2304.05203}{arXiv:2304.05203}.

\bibitem{Planck18VIII}
Planck Collaboration (N.~Aghanim \textit{et~al.}),
``Planck 2018 results. VIII. Gravitational lensing,''
\emph{Astron.\ Astrophys.} \textbf{641} (2020) A8,
\href{https://arxiv.org/abs/1807.06210}{arXiv:1807.06210}.

\bibitem{Brout22}
D.~Brout \textit{et~al.},
``The Pantheon+ analysis: cosmological constraints,''
\emph{Astrophys.\ J.} \textbf{938} (2022) 110,
\href{https://arxiv.org/abs/2202.04077}{arXiv:2202.04077}.

\bibitem{Rubin23}
D.~Rubin \textit{et~al.} (Union3),
``Union through UNITY: cosmology with 2,000 SNe using a unified
Bayesian framework,''
\href{https://arxiv.org/abs/2311.12098}{arXiv:2311.12098}.

\bibitem{DES5YR}
DES Collaboration,
``The Dark Energy Survey: cosmology results with $\sim 1500$
new high-redshift Type Ia supernovae using the full 5-year dataset,''
\href{https://arxiv.org/abs/2401.02929}{arXiv:2401.02929}.

\bibitem{DiValentino21}
E.~Di~Valentino \textit{et~al.},
``In the realm of the Hubble tension---a review of solutions,''
\emph{Class.\ Quant.\ Grav.} \textbf{38} (2021) 153001,
\href{https://arxiv.org/abs/2103.01183}{arXiv:2103.01183}.

\bibitem{Abdalla22}
E.~Abdalla \textit{et~al.},
``Cosmology intertwined: a review of the particle physics,
astrophysics, and cosmology associated with the cosmological tensions
and anomalies,''
\emph{JHEAp} \textbf{34} (2022) 49--211,
\href{https://arxiv.org/abs/2203.06142}{arXiv:2203.06142}.

\bibitem{MaggioreMancarella14}
M.~Maggiore and M.~Mancarella,
``Nonlocal gravity and dark energy,''
\emph{Phys.\ Rev.\ D} \textbf{90} (2014) 023005,
\href{https://arxiv.org/abs/1402.0448}{arXiv:1402.0448}.

\bibitem{Belgacem18}
E.~Belgacem, Y.~Dirian, S.~Foffa, and M.~Maggiore,
``Nonlocal gravity. Conceptual aspects and cosmological predictions,''
\emph{JCAP} \textbf{03} (2018) 002,
\href{https://arxiv.org/abs/1712.07066}{arXiv:1712.07066}.

\bibitem{Dirian14}
Y.~Dirian, S.~Foffa, N.~Khosravi, M.~Kunz, and M.~Maggiore,
``Cosmological perturbations and structure formation in nonlocal
infrared modifications of general relativity,''
\emph{JCAP} \textbf{06} (2014) 033,
\href{https://arxiv.org/abs/1403.6068}{arXiv:1403.6068}.

\bibitem{Maggiore16}
M.~Maggiore,
``Perturbative loop corrections and nonlocal gravity,''
\emph{Phys.\ Rev.\ D} \textbf{93} (2016) 063008,
\href{https://arxiv.org/abs/1603.01515}{arXiv:1603.01515}.

\bibitem{Belgacem19LLR}
E.~Belgacem, A.~Finke, A.~Frassino, and M.~Maggiore,
``Testing nonlocal gravity with Lunar Laser Ranging,''
\emph{JCAP} \textbf{02} (2019) 035,
\href{https://arxiv.org/abs/1812.11181}{arXiv:1812.11181}.

\bibitem{Belgacem19GW}
E.~Belgacem, Y.~Dirian, A.~Finke, S.~Foffa, and M.~Maggiore,
``Nonlocal gravity and gravitational-wave observations,''
\emph{JCAP} \textbf{11} (2019) 022,
\href{https://arxiv.org/abs/1907.02047}{arXiv:1907.02047}.

\bibitem{DeserWoodard07}
S.~Deser and R.~P.~Woodard,
``Nonlocal cosmology,''
\emph{Phys.\ Rev.\ Lett.} \textbf{99} (2007) 111301,
\href{https://arxiv.org/abs/0706.2151}{arXiv:0706.2151}.

\bibitem{NojiriOdintsov17}
S.~Nojiri, S.~D.~Odintsov, and V.~K.~Oikonomou,
``Modified gravity theories on a nutshell: inflation, bounce, and
late-time evolution,''
\emph{Phys.\ Rep.} \textbf{692} (2017) 1--104,
\href{https://arxiv.org/abs/1705.11098}{arXiv:1705.11098}.

\bibitem{Amendola17}
L.~Amendola, Y.~Dirian, H.~Nersisyan, and S.~Park,
``Observational constraints on nonlocal gravity with the $R\Box^{-2}R$
cosmological model,''
\emph{JCAP} \textbf{05} (2019) 052,
\href{https://arxiv.org/abs/1901.07832}{arXiv:1901.07832}.

\bibitem{CusinFoffaMaggiorMancarella16}
G.~Cusin, S.~Foffa, M.~Maggiore, and M.~Mancarella,
``Conformal symmetry and nonlinear extensions of nonlocal gravity,''
\emph{Phys.\ Rev.\ D} \textbf{93} (2016) 083008,
\href{https://arxiv.org/abs/1602.01078}{arXiv:1602.01078}.

\bibitem{BelgacemCusinFoffaMaggiore16}
E.~Belgacem, G.~Cusin, S.~Foffa, M.~Maggiore, and M.~Mancarella,
``Stability issues of nonlocal gravity during primordial inflation,''
\emph{Phys.\ Rev.\ D} \textbf{95} (2017) 104019,
\href{https://arxiv.org/abs/1610.05664}{arXiv:1610.05664}.

\bibitem{BelgacemDirianFoffaMaggiore18GW}
E.~Belgacem, Y.~Dirian, S.~Foffa, and M.~Maggiore,
``The gravitational-wave luminosity distance in modified gravity
theories,''
\emph{Phys.\ Rev.\ D} \textbf{97} (2018) 104066,
\href{https://arxiv.org/abs/1712.08108}{arXiv:1712.08108}.

\bibitem{NT4c}
D.~Alfyorov,
``NT-4c: FLRW reduction of SCT nonlocal field equations,''
SCT Theory internal document (2026).

\bibitem{NT4b}
D.~Alfyorov,
``NT-4b: Full nonlinear SCT spectral-action field equations on
arbitrary curved backgrounds,''
SCT Theory internal document (2026).

\bibitem{NT4a}
D.~Alfyorov,
``NT-4a: Linearized SCT spectral-action field equations on flat
background,''
SCT Theory internal document (2026).

\bibitem{INF1}
D.~Alfyorov,
``INF-1: Spectral inflation from the $R^2$ sector,''
SCT Theory internal document (2026).

\bibitem{PPN1}
D.~Alfyorov,
``PPN-1: Post-Newtonian solar system tests of SCT,''
SCT Theory internal document (2026).

\end{thebibliography}

\end{document}
